\section{What the incoming conductor audit establishes}
\label{sec:audit}
The supplied audit explicitly reports an incomplete reading of two source
excerpts. Its equations (1)--(17) are checked here against the complete
\emph{Weighted conductor forward} source, WCF1--27. It is evidence for
those finite calculations, not evidence that the unread programme lacks
their successors. In particular WCF19a--f already proves the full
determinant and inverse-exterior argument, WCF22 gives the actual lower
evaluation map, and WCF27 explicitly retains the error at the
$q\log a$ scale. The audit's full-determinant comparison is compatible
with these statements. We now prove a stronger finite bound using its
sharper weight estimate and the compression argument in WCF11a--12.

Fix $a=4\ell+1\ge9$, $q=(a+1)^2$, $q'=(a-7)^2$, the original
$0<\delta<1/2$, $\gamma>2$ and coefficients $a_{uw}$, $0\le u,w\le8$.
Put
\[
 E_A(z)=\sum_{u,w=0}^8 a_{uw}e^{[4+(2u-8)\delta+i(2w-8)\gamma]z},
 \quad v=\operatorname{ord}_0E_A,\quad \mu_v=E_A^{(v)}(0)\ne0,
 \quad A_{\rm abs}=\sum_{u,w}|a_{uw}|. \tag{AC1}
\]
The nonzero symbol is fixed as $a$ varies; $0\le v\le80$.
For $s\in\{1,a\}$, $b=s/2$, $N=D+v$, and $D\ge0$, retain
\[
 \begin{gathered}
 \rho_n^2=\frac{n!}{(b)_n},\quad
 w(t)=\tfrac12\log\frac{1-it}{1+it},\quad
 g(t)=t^{-v}e^{-4w(t)}E_A(w(t)),\\
 A_{D,s}=\operatorname{diag}(\rho_i)_{0\le i\le D}^{-1}
 T_D(g)\operatorname{diag}(\rho_{j+v})_{0\le j\le D}.
 \end{gathered}\tag{AC2}
\]
The logarithm is the branch vanishing at zero. The two original
polynomial spaces have the WCF2 measure
\[
 dm_s(y)=\frac{(2\pi)^{s/2}2^{s/2}}{4\pi\Gamma(s/2)}
 |\Gamma(s/4+iy/2)|^2dy
\]
and centres $a/2$ and $a/2-4$. No measure or leading symbol coefficient
is changed. Their orthonormal coordinates give (AC2); see WCF4--6,
whose derivation uses the Meixner--Pollaczek generating function
\cite[18.23.7]{DLMF18}.

\begin{proposition}[Sharpened original forward bound]
For $D\ge1$ define
\[
 \widehat C_N=e^{1+2\pi\gamma} A_{\rm abs}(2N+1)^{4\delta},\qquad
 \widehat B_{D,s}=
 \begin{cases}\widehat C_N(4N+1)^{1/4},&s=1,\\
 \widehat C_N,&s=a.
 \end{cases}\tag{AC3}
\]
Then $\|A_{D,s}\|\le\widehat B_{D,s}$. At $D=0$ its exact norm is
$|\mu_v|\rho_v/v!$. Every constant refers to the original symbol.
\end{proposition}
\begin{proof}
For $|t|=r<1$, the Cayley ratio lies in the right half-plane and has
modulus between $(1-r)/(1+r)$ and its reciprocal. Each exponent of that
ratio in $e^{-4w}E_A(w)$ has real part bounded by $4\delta$ and
imaginary part bounded by $4\gamma$. Hence
\[
 \max_{|t|=r}|g(t)|\le
 r^{-v} A_{\rm abs}e^{2\pi\gamma}
 ((1+r)/(1-r))^{4\delta}.\tag{AC4}
\]
For $b>1$, the exact coefficient norm is
\[
 \sum|c_n|^2\rho_n^2=(b-1)\int_0^1\int_0^{2\pi}
 |\sum c_n x^{n/2}e^{in\theta}|^2(1-x)^{b-2}
 \frac{d\theta}{2\pi}dx.
\]
Angular orthogonality and Euler's beta integral prove this identity
\cite[5.12.1]{DLMF5}. Multiplication by $g(rz)$ followed by projection
onto degrees at most $D$ therefore has norm at most $\max_{|t|=r}|g(t)|$.
Its transpose in the basis $z^n/\rho_n$ is
$\operatorname{diag}(\rho_i)^{-1}T_D(g(r\cdot))
\operatorname{diag}(\rho_i)$. Conjugating by
$\operatorname{diag}(1,r,\ldots,r^D)$ costs $r^{-D}$.
The remaining diagonal $\rho_{i+v}/\rho_i$ is a contraction. Thus
$\|A_{D,a}\|\le r^{-D}\max|g|$.

For $b=1/2$, use circle Parseval in place of the beta norm. The same
compression gives $\|T_D(g)\|\le r^{-D}\max|g|$ and the original
diagonal factors cost $\rho_N$, since $\rho_n$ increases from one.
We prove $\rho_n\le(4n+1)^{1/4}$ by induction. It holds at zero;
the squared induction step follows from
\[
 (4n+5)(n+\tfrac12)^2-(4n+1)(n+1)^2=\tfrac14>0.
\]
Choose $r=N/(N+1)$, which is allowed because $D\ge1$. Now
$r^{-D-v}=r^{-N}\le e$ and $(1+r)/(1-r)=2N+1$.
Combining these facts with (AC4) proves (AC3). For $D=0$, (AC2) is
the scalar $(-i)^v\mu_v\rho_v/v!$.
\end{proof}

\begin{proposition}[Determinant and original four-endpoint receiver]
Put $d=D+1$. The exact complex determinant and its two-source comparison are
\[
 \det A_{D,s}=((-i)^v\mu_v/v!)^d\prod_{j=0}^D\rho_{j+v}/\rho_j,
 \qquad
 |\det A_{D,s}|^2=|\mu_v/v!|^{2d}
 \prod_{j=0}^D\prod_{r=1}^v\frac{j+r}{j+b+r-1},\tag{AC5}
\]
\[
 0\le\log|\det A_{D,1}|-\log|\det A_{D,a}|
 \le\frac{v(a-1)}2 H_{D+1}.\tag{AC6}
\]
For the actual WCF lower evaluation matrices $E_i=\mathsf E_{-,D_i,s}$,
$N_i=(q-1,q,2q-1,2q)_i$, $D_i=N_i-v$, define
\[
 e_i=\log\frac{\det(E_i^*\overline A_{D_i,s}A_{D_i,s}^{\mathsf T}E_i)}
 {\det(E_i^*E_i)},\qquad \mathcal R e=-e_0-e_1+e_2+e_3.
\]
Set $M_i=\max(1,\widehat B_{D_i,s})$ and
\[
 J_i=(D_i+1)\log^+|v!/\mu_v|
       +\tfrac12(s/2-1)_+vH_{D_i+1},\quad m=16a-48-v.
\]
Then the original conductor correction satisfies
\[
 \mathcal R e\le
 2m\log M_0+2(m+1)\log M_1+2J_0+2J_1
 +2q'\log M_2+2q'\log M_3.\tag{AC7}
\]
For the fixed actual symbol, this improves WCF26a to
\[
 \mathcal R e\le
 \begin{cases}
 (32\delta+2)q\log a+O_{\rm actual}(q),&s=1,\\
 32\delta q\log a+O_{\rm actual}(q),&s=a.
 \end{cases}\tag{AC8}
\]
This is an upper bound, not an assigned endpoint value.
\end{proposition}
\begin{proof}
Triangularity proves (AC5). Dividing its two squared moduli gives
one half the sum of
$\log(1+(a-1)/(2(j+r-1/2)))$. Use $\log(1+x)\le x$ and
$j+r-1/2\ge(j+1)/2$ to prove (AC6). The gamma-product version in
the audit follows by multiplying the consecutive factors in $j$.
For fixed $v$ and $D=O(a^2)$, (AC6) divided by $q$ tends to zero.
This comparison is at the same dimension and concerns the full matrix.

For any invertible $d\times d$ matrix $B$, its singular values give
$\|\bigwedge^rB^{-1}\|=\|\bigwedge^{d-r}B\|/|\det B|$.
Formula (AC5) and $\log(1+x)\le x$ give
$|\det A_{D_i,s}|^{-1}\le e^{J_i}$. Consequently the determinant
of the Gram of any $q'$ independent vectors after applying
$A_{D_i,s}^{\mathsf T}$ lies between its original determinant times
$e^{-2((D_i+1-q')\log M_i+J_i)}$ and times $M_i^{2q'}$.
This follows by applying the exterior map, then its inverse, to their
full wedge. WCF22 identifies these vectors with the actual conductor
evaluation columns; no arbitrary endpoint substitutes for them.
At the low endpoints $D_0+1-q'=m$ and $D_1+1-q'=m+1$.
Insert the lower bounds at those two endpoints and upper bounds at the
two high endpoints to obtain (AC7), with its stated signs.

For the high endpoints $N_i=2(a+1)^2+O(1)$, (AC3) gives
$\log M_i\le8\delta\log a+O_{\rm actual}(1)$ at $s=a$, and
$(8\delta+1/2)\log a+O_{\rm actual}(1)$ at $s=1$.
The low-dimensional first two terms in (AC7) are $O(a\log a)$;
$J_0+J_1=O_{\rm actual}(q+a\log a)$. Since $q'\le q$, (AC8)
follows. The original WCF27 complete action receives $\mathcal R e$
with coefficient one; replace its corresponding upper allowance by
the minimum of the old allowance and (AC7). All other action terms remain.
\end{proof}

The audit's trace-log remainder follows by diagonalizing its Hermitian
$H$ and bounding the scalar series remainder by
$|\lambda|^{M+1}/((M+1)(1-\theta))$. Summation yields both its
rank and trace-norm versions. Its Vandermonde estimate follows from
$\sum_n|[z^n]\prod_{k\ne j}(z-\beta_k)|\le(1+B)^{80}$ and
$\prod_{k\ne j}|\beta_j-\beta_k|\ge(2\delta)^{80}$; it bounds
the largest of the first 81 moments, not necessarily the first nonzero
one. These are valid checks, with the stated scopes. Neither implies
individual projected-current signs or absolute native allocations.
